\documentclass[10pt,a4paper]{amsart}
\usepackage[margin=19mm]{geometry}
\usepackage{amsmath,amssymb,mathtools}
\usepackage[scr=rsfs,cal=euler]{mathalfa}
\usepackage{xfrac}
\usepackage[unicode,hidelinks,bookmarks=false]{hyperref}
\newcommand{\ie}{i.e.\ }
\newcommand{\cf}{cf.\ }
\newcommand{\resp}{respectively\ }
\newcommand{\Subsection}{\subsection}
\providecommand{\nobreakdash}{\nobreak\hbox{-}}
\setlength{\emergencystretch}{2em}
\multlinegap0pt
\usepackage{bm}
\usepackage{bbm}
\usepackage{dsfont}
\usepackage{xspace}
\usepackage{cancel}
\usepackage{mathtools}
\usepackage{amsthm}
\newtheorem{thm}{Theorem}
\newcommand{\Q}{\mathbb{Q}}
\title[Beyond the unique pair: a two-parameter family of rational t]{Beyond the unique pair: a two-parameter family of rational triangles with equal perimeter and area (Theorem 9.1)}
\author{Jesse Allen, Olivier Beaudry-Ogden, Th\'eophane Bouchard, Mat\'eo Fr\'ely, Matilde Lal\'in, Abel-Jimmy Oyono-Montoki,, Berend Ringeling}
\date{Source: arXiv:2609.28020v1 (CC BY 4.0)}
\begin{document}
\maketitle

\noindent\emph{Source and licence.} Beyond the unique pair: a two-parameter family of rational triangles with equal perimeter and area. Version arXiv:2609.28020v1 (CC BY 4.0), distributed under the Creative Commons Attribution 4.0 International licence, as recorded in the arXiv licence metadata for this exact version and shown on the abstract page https://arxiv.org/abs/2609.28020v1. Full rights notice: licenses/arxiv-2609.28020-CC-BY-4.0.txt.\par\smallskip

\noindent\textit{Editorial scope.} Counted unit: the Thm printed as Theorem 9.1 of the article (source0.tex lines 2581--2598, source label \texttt{thm:four-pair-rational}), delivered together with the article's own complete original proof of it (source0.tex lines 2600--2720, one \texttt{proof} environment of 121 lines). No mathematics is written, repaired or re-derived: statement and proof are the article's own text line for line. Required context retained uncounted: none. Every definition, display and label of those lines is retained unchanged. Omitted from this package and not redistributed: 1--2580, 2599--2599, 2721--5073. Nothing inside a retained excerpt is omitted or altered: every retained body is delivered exactly as the article's lines are written, so no omission receipt of any kind accompanies this package. Citations: the retained excerpts carry no citation command at all, so no bibliography entry of the article is transcribed and no reference is redistributed. Reference adaptation: none was needed and none was made; every reference inside the delivered unit resolves inside it, so no unresolved reference or citation is produced. Printed slips kept literal: no printed word, symbol, hypothesis or number of the counted statement or of its proof was altered. Layout and class: the article's own class is \texttt{amsart} with packages including amsmath, amsthm, amssymb, booktabs, array, multirow, float, graphicx, xcolor, longtable, geometry, hyperref; the wrapper is the lane's pinned \texttt{amsart} editorial wrapper at 10pt on A4, which restates the theorem-like environments the retained text needs and adds the article's own macro definitions that the retained text needs (\texttt{\textbackslash Q}), with extra wrapper packages (bm, bbm, dsfont, xspace, cancel, mathtools). The delivered numbering is the unit's own. Assets: no retained span contains a figure environment, a graphics-inclusion command, a PDF-page inclusion, a table or a TikZ picture; no image file is copied into this package, the package declares no asset dependency and no image file is redistributed with it. Licence: version arXiv:2609.28020v1 of this article is distributed under the Creative Commons Attribution 4.0 International licence, as recorded in the arXiv licence metadata for this exact version and shown on the abstract page https://arxiv.org/abs/2609.28020v1. A full rights notice is delivered at licenses/arxiv-2609.28020-CC-BY-4.0.txt. Characters: every retained excerpt is pure ASCII, so the delivered engine, XeLaTeX with its default Latin Modern text font, reports no missing character.
\begin{thm}\label{thm:four-pair-rational}
For every
$\lambda\in\Q$ and $\frac13\leq\lambda\leq\frac37$,
the value
$\rho=\rho_4(\lambda)$
admits at least four distinct genuine pairs, namely
\[
\bigl(T_0,I(z_-)\bigr),
\qquad
\bigl(T_0,I(z_+)\bigr),
\qquad
\bigl(T_1,I(\alpha)\bigr),
\qquad
\bigl(T_1,I(\beta)\bigr).
\]
In particular, there are infinitely many rational values of $\rho$
carrying at least four distinct genuine pairs.
\end{thm}

\begin{proof}
All six triangles have perimeter $2$. We first check that $T_0$ and
$T_1$ have an angle $\theta$ with
$\cos(\theta)=\rho_4(\lambda)$.
In both cases we take the distinguished angle to be the angle between
the last two sides. Using
$\delta^2+3r^2=1,$
the law of cosines applied to $T_0$ gives
\[
\frac{(2r^2+\delta)^2+(2r^2-\delta)^2-(2-4r^2)^2}
     {2(2r^2+\delta)(2r^2-\delta)}
=
\frac{1-r^2}{1+r^2},
\]
and similarly to $T_1$,
\[
\frac{\left(\frac{2+\delta}{3}\right)^2
      +\left(\frac{2-\delta}{3}\right)^2
      -\left(\frac23\right)^2}
     {2\left(\frac{2+\delta}{3}\right)
        \left(\frac{2-\delta}{3}\right)}
=
\frac{1-r^2}{1+r^2}.
\]
Thus both $T_0$ and $T_1$ are rational $\theta$-triangles for
$\rho=\rho_4(\lambda)$.

For an isosceles triangle $I(z)$ with $0<z<1$, its area is
$\operatorname{Area}(I(z)) = \frac{z(1-z^2)}{2}$.
Since
$z_\pm=r\pm\delta$
and $\delta^2=1-3r^2$, we have
$1-z_\pm^2 = 2r(r\mp\delta)$.
Therefore,
\[
\operatorname{Area}(I(z_-))
=
\operatorname{Area}(I(z_+))
=
r(r^2-\delta^2)
=
r(4r^2-1).
\]

On the other hand, Heron's formula, together with
$\delta^2+3r^2=1$, gives
$\operatorname{Area}(T_0)=r(4r^2-1)$ and similarly, $\operatorname{Area}(T_1) = \frac r3$.


A direct calculation also gives
$r = \frac32\alpha(1-\alpha^2) = \frac32\beta(1-\beta^2)$.
Therefore
$\operatorname{Area}(I(\alpha)) = \operatorname{Area}(I(\beta)) = \frac r3$.

Thus the four displayed pairs have equal perimeter and equal area.

It remains to check that they are genuine and distinct. On the interval
$\frac13\leq\lambda\leq\frac37$
one has
\[
\frac12<r<\frac1{\sqrt3},
\qquad
0<\delta<\frac12,
\qquad
0<z_-<z_+<1,
\]
and
\[
\frac37\leq \alpha\leq\frac7{13}
<
\frac8{13}\leq \beta\leq\frac57.
\]
Hence all four isosceles triangles are non-degenerate.

For $T_0$, the preceding bounds and the identity
$\delta^2+3r^2=1$ give
\[
0<2r^2-\delta
<
2-4r^2
<
2r^2+\delta
<1,
\]
since $1-(2r^2+\delta)
=\frac{(1-2\delta)(1-\delta)}{3}>0.$ Since the three sides lie in $(0,1)$ and the perimeter is 2,  $T_0$ is a non-degenerate scalene triangle. Likewise,
\[
0<
\frac{2-\delta}{3}
<
\frac23
<
\frac{2+\delta}{3}
<1.
\]
Since the perimeter is 2,  $T_1$ is also non-degenerate and scalene.

The inequalities $z_-\ne z_+$ and $\alpha\ne \beta$ show that each of
$T_0$ and $T_1$ has two distinct isosceles partners. Moreover,
$T_0$ and $T_1$ are not congruent, as  the side opposite the distinguished
angle has length
$2-4r^2>\frac23$
for $T_0$, while it has length $\frac23$ for $T_1$. Hence the four
displayed pairs are distinct and genuine.

Finally, for $\frac13\leq\lambda\leq\frac37$,
\[
r'(\lambda)
=
-\frac{
6(\lambda^2-3)(\lambda^2-6\lambda-3)
(\lambda^2+6\lambda-3)}
{(\lambda^2+3)^4}
>0.
\]Thus $r(\lambda)$ is strictly
increasing, and therefore
$\rho_4(\lambda)=\frac{1-r(\lambda)^2}{1+r(\lambda)^2}$
is strictly decreasing. Since there are infinitely many rational
$\lambda$ in the stated interval, this gives infinitely many distinct
rational values of $\rho$ carrying at least four genuine pairs.
\end{proof}

\end{document}
